\begin{ex}\label{es:1}
Fissiamo l'alfabeto binario $A=\{0,1\}$, dimostrare che l'insieme
\[
X = \{ w0 \mid w\in A^* \}
\]
è decidibile per automa a stati finiti.


\begin{sol} Si richiede di mostrare un automa il cui algoritmo associato accetta le parole che terminano con uno zero. È dunque sufficiente memorizzare nello stato l'ultimo simbolo letto e alla fine della lettura della parola decidere in base allo stato: poniamo $Q=\{q_0,q_1\}$, e definiamo
\[
T(q,a) :=
\begin{cases}
q_0 & \text{se } a=0, \\
q_1 & \text{se } a=1,
\end{cases}
\]
prendiamo $q_0$ come stato iniziale e $F=\{q_0\}$ come insieme degli stati finali accettanti.
.
Sia
\[
X=\{w0\mid w\in A^*\}
\]
l'insieme delle parole sull'alfabeto binario che terminano con il simbolo \(0\).

Costruiamo un automa finito che decide \(X\). Consideriamo l'insieme degli stati
\[
Q=\{q_0,q_1\},
\]
dove intuitivamente:
\begin{itemize}
    \item \(q_0\) rappresenta la situazione in cui l'ultimo simbolo letto è \(0\);
    \item \(q_1\) rappresenta la situazione in cui l'ultimo simbolo letto è \(1\).
\end{itemize}

Definiamo la funzione di transizione
\[
T:Q\times A\to Q
\]
come
\[
T(q,a)=
\begin{cases}
q_0 & \text{se } a=0,\\
q_1 & \text{se } a=1.
\end{cases}
\]

Scegliamo come stato iniziale
\[
q_{\mathrm{init}}=q_0
\]
e come insieme degli stati accettanti
\[
F=\{q_0\}.
\]

Consideriamo ora l'algoritmo associato all'automa.
Sia
\[
w=a_1a_2\cdots a_n\in A^*.
\]
La configurazione iniziale è
\[
\epsilon(w)=(w,1,q_0).
\]

Dopo \(k\) transizioni, per induzione su \(k\), la configurazione raggiunta è
\[
(w,k+1,q_k),
\]
dove
\[
q_k=
\begin{cases}
q_0 & \text{se } a_k=0,\\
q_1 & \text{se } a_k=1.
\end{cases}
\]

Infatti, ad ogni passo la funzione di transizione sostituisce lo stato corrente con lo stato corrispondente all'ultimo carattere letto.

Dopo aver letto tutta la parola \(w\), la computazione termina nella configurazione
\[
(w,n+1,q_n).
\]
La funzione di uscita restituisce quindi
\[
\delta(w,n+1,q_n)=
\begin{cases}
1 & \text{se } q_n\in F,\\
0 & \text{se } q_n\notin F.
\end{cases}
\]

Poiché
\[
F=\{q_0\},
\]
si ha
\[
\delta(w,n+1,q_n)=
\begin{cases}
1 & \text{se l'ultimo simbolo di } w \text{ è }0,\\
0 & \text{se l'ultimo simbolo di } w \text{ è }1.
\end{cases}
\]

Quindi
\[
f_{\mathcal A}(w)=
\begin{cases}
1 & \text{se } w\in X,\\
0 & \text{se } w\notin X.
\end{cases}
\]

Pertanto
\[
f_{\mathcal A}=\mathbf 1_X
\]
e quindi l'algoritmo associato all'automa decide \(X\).

Concludiamo che
\[
X=\{w0\mid w\in A^*\}
\]
è DFA-decidibile.
\end{sol}
\end{ex}

\newpage

\begin{ex}\label{es:3}
Si consideri l'alfabeto $A=\{0,1\}$ e sia $B=\{1\}$ la restrizione all'insieme che contiene il solo simbolo $1$; si considerino l'insieme delle parole di alfabeto $A$ di lunghezza pari:
\[
X_1 = \{ w \in A^* \mid |w| = 2n,\ n\geq 0 \}
\]
e l'insieme delle parole di alfabeto $A$ che contengono un numero pari di $1$:
\[
X_2 = \{ w \in A^* \mid |w_{|B}| = 2k,\ k\geq 0 \}
\]
qui si è utilizzata la notazione $w_{|B}$ per indicare la parola $w$ 
ristretta ai soli elementi di alfabeto $B$.
Costruire l'automa che decide $X_1+X_2$.

\begin{sol}
Utilizzando il teorema 
\ref{prop:8},
basta esibire gli automi che decidono $X_1$ e $X_2$.
L'automa che decide $X_1$ è quello dell'esempio 
\ref{ex: parole lunghezze pari} (il grafo corrsipondente e' nella figura 
\ref{figura esempio automa per parole pari}).
Definiamo ora l'automa che decide $X_2$.
Definiamo $Q_2=\{q_0, q_1 \}$ dove $q_0$ e/' sia stato iniziale che stato accettante, mentre $q_1$ e/' uno stato non accettante.
E definiamo la funzione di transizione $T_2$ come segue:
\[
\begin{matrix}  
T_2: & Q_2 \times A & \to  & Q_2
\\
& (q_0,0) & \mapsto & q_0\\
& (q_0,1) & \mapsto & q_1\\
& (q_1,0) & \mapsto & q_1\\
& (q_1,1) & \mapsto & q_0   
\end{matrix}
\]
In questo modo, l'automa che decide $X_2$ 
accetta le parole che contengono un numero pari di $1$
(si guardi il grafico dell'automa nella figura \ref{fig:parita}).
A questo punto,
utilizzando la costruzione 
effettuata
nella dimostrazione 
del teorema \ref{prop:8}
otteniamo che l'automache decide 
$X_1 + X_2$
e/' l'automa
\[
\begin{cases}
    Q:= Q_1 \times Q_2,\\
    q_0 := (q_0',q_0''),\\
    F := (F_1 \times Q_2) \cup (Q_1 \times F_2),\\
    T((q_1,q_2),a) = (T_1(q_1,a), T_2(q_2,a))
\end{cases}
\]
dove con un apice indichiamo gli stati di dell'automa che decide 
$X_1$ e con due apici gli stati dell'automa che decide $X_2$.
Per completezza,
presentiamo il grafo dell'automa che decide $X_1 + X_2$
(che puo/' esssere disegnato a partire dall'automa appena definito
oppure usando i grafi degli automi che decidono $X_1$ e $X_2$).
\begin{center}
\begin{tikzpicture}[shorten >=1pt, node distance=4cm, on grid, auto]
   % Stati disposti a quadrato seguendo il ciclo delle transizioni
   \node[state, initial, accepting] (q00) {$(q_0', q_0'')$};
   \node[state, accepting]          (q10) [right=of q00] {$(q_1', q_0'')$};
   \node[state, accepting]          (q01) [below=of q10] {$(q_0', q_1'')$};
   \node[state]                     (q11) [left=of q01]  {$(q_1', q_1'')$};

   % Transizioni senza intersezioni
   \path[->]
    % Bordo superiore (Transizioni con 0)
    (q00) edge [bend left=15] node {0} (q10)
    (q10) edge [bend left=15] node {0} (q00)
    
    % Bordo destro (Transizioni con 1)
    (q10) edge [bend left=15] node {1} (q01)
    (q01) edge [bend left=15] node {1} (q10)
    
    % Bordo inferiore (Transizioni con 0)
    (q01) edge [bend left=15] node {0} (q11)
    (q11) edge [bend left=15] node {0} (q01)
    
    % Bordo sinistro (Transizioni con 1)
    (q11) edge [bend left=15] node {1} (q00)
    (q00) edge [bend left=15] node {1} (q11);
\end{tikzpicture}
\end{center}
\end{sol}
\end{ex}